\chapter{Automated Numerical Framework}\label{chap:numerical_model}

\section{Framework Overview}

To complement the experimental campaign presented in
Chapter~\ref{chap:experimental_campaign}, a numerical simulation framework was
developed in Abaqus/Explicit. Its purpose is to create, run, and post-process
the specimen and punch configurations considered in Task~I in a reproducible
way. The framework is not written for one single model, but for repeated
simulation campaigns in which the test type, geometry, mesh density, material
orientation, loading, and post-processing settings can be changed from a common
interface.

The main design goal was automation. Starting from a central configuration
file, the pipeline builds the finite element model, writes the Abaqus input
deck, submits the simulation, extracts the relevant quantities from the output
database, and stores the processed results in a common format. This makes it
possible to compare Nakazima, Marciniak, and Punch-in-Punch simulations using
the same numerical and post-processing workflow.

Special attention was paid to the mesh. Since the numerical results should be
compared across different specimen and punch configurations, the element size
must be controlled systematically. The framework therefore includes
user-defined mesh-size modification routines, which allow mesh-refinement
studies to be generated automatically instead of manually rebuilding each
model.

\section{Framework Architecture}
\subsection{Central Configuration}

The framework is controlled through a central configuration file, which acts as
the user interface of the numerical model. It defines the selected test
configuration, specimen dimensions, punch setup, mesh controls, material
orientation, contact parameters, loading conditions, solver settings, and
post-processing options. Parameters that are expected to change between
campaigns are exposed there, while the implementation of the model-building
steps remains fixed.

This separation reduces the risk of introducing unintended differences between
simulation cases. For example, a change in sheet thickness, punch diameter, or
mesh refinement factor does not require editing the Abaqus model manually. The
complete list of available configuration groups is given in
Appendix~\ref{app:configuration_parameters}.

\subsection{Automated Workflow}

The model generation is organised as a sequence of Python modules. The main
script first reads the configuration and then creates the specimen, tools,
assembly, material definition, analysis step, contact interactions, boundary
conditions, output requests, and job definition. Each simulation therefore
passes through the same sequence of operations, independent of the selected
test type.

After the model has been generated, the input file and the user material
subroutine are placed in the simulation directory. Once the Abaqus/Explicit job
has finished, the same post-processing script reads the output database (ODB)
and writes the force-displacement data, energy histories, strain paths, and
forming-limit points. The resulting file structure is kept identical for all
test configurations so that the results can be aggregated without additional
manual work.

\subsection{Integration with the ETH Euler Cluster}

The simulations are executed on the ETH Euler cluster. The local project files
are copied to the cluster, where Abaqus/CAE is used in no-GUI mode to build the
model and write the input deck. The solver job is then submitted through SLURM,
and the ODB post-processing is performed after the analysis has completed.

Only the compact result files, figures, and movies are copied back to the
project directory. The large ODB files remain in scratch storage. This keeps
the workflow practical for complete specimen-width campaigns and for mesh or
mass-scaling studies, where many independent simulations have to be submitted.

\section{Automated Model Generation}

\subsection{Test Configurations}

Three test configurations are supported. The Nakazima model represents the
standard hemispherical-punch test and is used as the main reference
configuration. The Marciniak model uses a flat punch and provides a second
forming setup with a different strain-path character. The Punch-in-Punch model
contains an outer and an inner punch and is used to study a two-stage forming
process.

All three configurations are generated by the same framework. Differences such
as the number of punches, tool geometry, contact pairs, and loading sequence
are selected through the configuration file. This common structure is important
because it allows the forming-limit points from the different tests to be
compared in the same post-processing pipeline.

\subsection{Geometry, Contact and Boundary Conditions}

To reduce computational cost, one quarter of the specimen is modelled. Symmetry
boundary conditions are applied on the two cut planes, while the tools are
treated as rigid bodies. The punch motion is prescribed through its reference
point, and the die and blank holder are fixed. The outer blank edge is
constrained to represent the clamping action of the experimental setup.

Tool--blank contact is defined with hard normal contact and Coulomb friction.
For the standard Nakazima and Marciniak models the same contact formulation is
used, while the Punch-in-Punch configuration activates the additional contact
pairs required by the second punch. In this way, the contact definition remains
consistent but can still represent the different experimental setups.

\subsection{Parameterised Mesh Generation}

The deformable sheet is discretised with solid brick elements. A solid
formulation is used because the material model includes through-thickness
deformation, damage evolution, and element deletion. The mesh is generated or
adapted automatically from the selected specimen geometry.

The mesh routines are designed to keep the element size comparable between
different cases. A mesh refinement factor can be applied to scale the target
element size, which makes it possible to run a mesh-size study from the same
model definition. This is especially important for forming-limit simulations,
because the extracted failure strain can be sensitive to the spatial resolution
in the necking region. The detailed mesh-study grid is therefore reported with
the results rather than fixed in the model description.

\section{Numerical Formulation}

\subsection{Material Model}

The sheet material is described by a user material subroutine. The model
combines anisotropic plasticity, mixed hardening, and a Hosford--Coulomb damage
criterion. The material orientation can be changed in the configuration file,
which allows rolling-direction and transverse-direction cases to be simulated
with the same model setup.

Element deletion is activated through the damage variables of the material
model. This makes it possible to identify the fracture location in the
post-processing step and to distinguish valid dome fractures from failures that
occur outside the evaluation zone.

\subsection{Explicit Solution Strategy}

All simulations are run with Abaqus/Explicit. Although an explicit solver is
used, the forming process is intended to represent a quasi-static experiment.
For this reason, the punch velocity and mass-scaling settings are controlled
carefully, and each run is checked using the ratio of kinetic energy to
internal energy.

The energy history is exported for every simulation. A low kinetic-to-internal
energy ratio indicates that inertial effects remain small and that the
simulation can be interpreted as quasi-static. This check is used together with
the stability of the extracted forming-limit point when selecting suitable
solver and mesh settings.

\section{Automated Post-Processing and Forming-Limit Extraction}

After each simulation, the ODB is processed automatically. The post-processing
script extracts the punch force-displacement curve, global energy histories,
local strain paths, element status information, and the quantities required for
forming-limit identification. These data are written to CSV files with the same
schema for every test configuration.

The fracture endpoint is detected primarily from the punch force history and
is checked against the element-deletion information. The critical strain is
then extracted from the dome region. In addition to the fracture limit, the
pipeline can write necking indicators based on the Volk--Hora method and on
the damage evolution of the material model when the required data are
available.

This automated post-processing step closes the numerical workflow. It converts
the raw finite element output into the same type of forming-limit data used in
the experimental evaluation, allowing the numerical and experimental results
to be compared in Chapter~\ref{chap:results}.
